\input{\string~/Documents/School/"UC Irvine"/ta_template2.0.tex}
\rh{2B}{5.4}
\chead{\today}
\graphicspath{ {\string~/Desktop} }
\usetikzlibrary{arrows}
%============ANSWER KEY TOGGLE===========
\newcommand{\ans}[1]{ \noindent {\color{red} \textbf{Answer:} #1}}
%\renewcommand{\ans}[1]{\vspace{3in}}
%==========================================
\begin{document}
\begin{center}
\textbf{Indefinite Integrals}
\end{center}

\begin{aun}
The big idea in this section is just to remember why we have to use $+C$ when we are taking antiderivatives: both $x^2 + 3$ and $x^2 + 5$ have the same derivative, even though they're different functions. In general, derivatives tend to lose some information, and the $C$ accommodates for all possibilities of that lost constant. 
\end{aun}


\begin{exx}
Complete the following table of antiderivatives
\[
\begin{array}{|c|c|c|c|}
\hline
\text{Function} & \text{Antiderivative} & \text{Function} & \text{Antiderivative}\\ \hline
x^n & & \frac{1}{x} &  \\
C & & e^x & \\
\cos(x) & & \sec^2(x) & \\
\csc(x)\cot(x) & & \frac{1}{\sqrt{1-x^2}} & \\ 
\frac{1}{x\sqrt{x^2 - 1}} & & \frac{-1}{\sqrt{1-x^2}} & \\\hline
\end{array}
\]
\end{exx}

\ans{
\[
\begin{array}{|c|c|c|c|}
\hline
\text{Function} & \text{Antiderivative} & \text{Function} & \text{Antiderivative}\\ \hline
x^n & \frac{x^{n+1}}{n+1} +C & \frac{1}{x} & \ln|x| + C \\
C & Cx + D & e^x & e^x + C\\
\cos(x) & \sin(x) + C& \sec^2(x) & \tan(x) + C \\
\csc(x)\cot(x) & -\csc(x) + C & \frac{1}{\sqrt{1-x^2}} & \arcsin(x) + C \\
\frac{1}{x\sqrt{x^2-1}} & \sec^{-1}(x) + C & \frac{-1}{\sqrt{1-x^2}} & \arccos (x) + C \\\hline
\end{array}
\]
}

\problembox{Compute $f(x)$ when:
\[
	f''(x) = e^x -2\sin (x), \quad f(0)=3 \quad \text{ and} \quad f(\pi/2) = 0
\]}

\ans{
We integrate:
\[
\int e^x -2 \sin(x) dx = e^x + 2\cos(x) + C = f'(x)
\]
...and once more
\[
\int e^x + 2\cos(x) + C dx = e^x + 2\sin(x) + Cx + D = f(x)
\]
Now, using the $f(0)= 3$, we test:
\[
f(0) = e^0 + 2\sin(0) + C\cdot 0 + D = 1 + D = 3
\]
implying $D = 2$. Similarly, we use $f(\pi/2) = 0$ to find:
\[
f(\pi/2) = e^{\pi/2} + 2\sin(\pi/2) + C \cdot \pi/2 + 2 = 0
\]
implying that 
\[
C = -\frac{2}{\pi} \lrp{e^{\pi/2} +4}
\]
This gives us the final answer:
\[
\boxed{f(x) = e^x + 2\sin(x) + -\lrp{\frac{2}{\pi} \lrp{e^{\pi/2} +4}}x + 2}
\]
}

\problembox{Compute the indefinite integral:
\[
\int \frac{\cos x}{\sin^2 x} dx
\]}

\ans{We know that $\frac{\cos}{\sin} = \cot$, and what is left is a factor of $\frac1\sin = \csc$:
\[
\int \csc x \cot x dx = -\csc x + C
\]}

\problembox{Evaluate the following definite integral:
\[
\int_0^2 x \lrp{\sqrt[3]{x} + \sqrt[4]{x}} dx
\]}

\ans{We compute the general antiderivative:
\[
\int x \lrp{\sqrt[3]{x} + \sqrt[4]{x}} dx = \int x^{4/3} + x^{5/4}dx
\]
When we take the ``reverse power rule'' we have:
\[
\frac37 x^{7/3} + \frac49 x^{9/4} + C
\]
and plug in:
\[
[\frac37 x^{7/3} + \frac49 x^{9/4} + C]_0^2 = 2^{7/3} + 2^{9/4} 
\]
}

\end{document}